\chapter{结构化初始化与调度算法改进}

\label{sec:algorithm-cn}

针对配对、腔室周期和共享资源次序共同决定的组合搜索空间，本文提出基于强化学习的双源混流组合设备调度算
法（reinforcement learning for dual-source mixed-flow cluster tool scheduling，RL-DMCT）。算
法输入为Pjob、LP库存、实际编号映射及设备参数，输出为具体晶圆动作表和产品完工时间等调度指标。

RL-DMCT通过改进初始解和候选割选择提高限时求解效率。初始解策略先保留配对、腔室分配和加工周期，再逐步放开
资源使用次序进行优化，为主问题提供较好的可行解。选割时，将候选割中的变量角色特征输入学习策略，再根据
策略排序、变量角色覆盖程度和候选割相似度进行重排。主问题求解后，继续优化动作时刻，并解码得到每片晶圆的排程。

\section{基于模型结构的初始解策略}
\label{sec:pjob-encoding-cn}

可行排程需要同时满足产品配对、满片加工、连续处理、清洁与辅助片借还约束。在完整模型中直接搜索时，找到首个可行解可能占用较多计算时间，而提高松弛下界也不能直接得到可执行排程。因此，先根据工艺要求构造配对、分腔和周期安排，再逐步调整资源使用次序，以减少初期需要同时调整的离散决策，并为后续联合搜索提供可行上界。

初始解按任务分配、PW对组成和加工周期依次构造：先将产品划分为连续区段并分配到允许腔室，再按工艺
模式组成PW对，以辅助片补齐空位；随后安排生产周期、$2\times2$连续处理链及清洁边界，沿供片、加工和回片
路径，结合操作时长、机械手移动与资源释放时刻依次确定各事件的开始时间。所得排程经出片顺序、资源占用和
驻留条件检查后记为$\mathcal S^0$。将其中的产品对、辅助身份、分腔、周期占位及动作时刻对应到模型变量，
并补齐资源次序与逻辑辅助量，得到完整初始赋值$\xi^0=\mathsf{Enc}(\mathcal S^0)$；其中$\mathsf{Enc}$
表示物理排程到模型赋值的映射。

采用结构保持二元骨架（structure-preserving binary skeleton，SPBS），逐层放开部分离散变量，以缩短$C_{\max}$为目标改进初始排程。第一层固定主要离散选择，调整连续时刻和取放同步选择；随后依次开放VTR
次序、装载锁次序和其余资源次序；最后开放周期占位、前后侧及手分配，保留PW对组成、辅助身份、分腔、库存指派
和周期类型。令$t\in\{1,\ldots,5\}$为层号，$\mathcal F_t$为该层固定的离散分量索引集合，以该层起点$\widehat\xi$
为参考，临时约束为：
\begin{equation}
\xi_d=\widehat\xi_d\quad(d\in\mathcal F_t),\qquad
\mathcal F_1\supseteq\cdots\supseteq\mathcal F_5.
\label{eq:spbs-fix-cn}
\end{equation}
每层至多使用初始解改善预算的五分之一，并将当前最好可行赋值传给下一层。该策略先优化给定配对、分腔和周期安排下的动作时刻，再调整资源使用次序，使搬运与加工能够重叠执行。改善后的赋值经完整验证后作
为第一阶段的起点，其产品完工时间提供可行上界，进入第一阶段时解除全部临时固定关系。已有可行赋值可直接作
为SPBS起点，构造与改善成本同主求解时间分别记录。


\section{结构特征感知重排与强化学习选割}
\label{sec:hem-extension-cn}

候选割的分离强度不同，涉及的变量和约束方向也可能相近。当单轮选割数量有限时，只按各割的质量排序可能选入多个相似割。因此，RL-DMCT先提取候选割的通用质量特征和变量角色特征，由强化学习策略确定选择数量和候选顺序，再根据已选割的变量类型及候选割相似度进行重排。学习排序采用HEM的层次数量控制与有序子集生成框架\citep{wang2023hem}，变量角色特征同时用于策略输入和重排评分。

\paragraph{变量角色特征提取}
每个候选割首先提取13维通用特征：目标平行度、有效性、支持比例、整数支持、归一化违约程度，以及割系数和目标系数
各自的均值、最大值、最小值与标准差。进一步按变量类型和目标角色追加10维结构特征。设$\mathcal V_a$为割$a$的非
零系数变量集合，$\alpha_{av}$为变量$v$在该割中的系数；$\mathcal V_g$依次表示二元、一般整数、连续目标、连续
辅助及其他五类变量，$g\in\{1,\ldots,5\}$。连续目标变量指目标系数非零的连续变量。各角色的系数占比为：
\begin{equation}
p_{ag}=\frac{\sum_{v\in\mathcal V_a\cap\mathcal V_g}|\alpha_{av}|}
{\sum_{v\in\mathcal V_a}|\alpha_{av}|}.
\label{eq:role-profile-cn}
\end{equation}
该统计计算各类变量的绝对系数和占总绝对系数和的比例，消除整行正比例缩放的影响，且不受同类变量编号顺序的影响。对于
本文模型，二元角色包含配对、分腔与资源排序选择，连续角色包含目标量和动作时刻，角色划分直接读取MILP变量属性。

在五类占比之外，提取有限上下界变量的系数占比、正系数占比、离散与连续角色的耦合量$\chi_a$、主导已识别角色占
比及归一化角色熵$\mathcal H_a$，组成23维输入；熵中的零项按零处理。其中：
\begin{equation}
\begin{aligned}
\chi_a&=4(p_{a1}+p_{a2})(p_{a3}+p_{a4}),\\
\mathcal H_a&=-\frac{\sum_{g=1}^{5}p_{ag}\log p_{ag}}{\log5}.
\end{aligned}
\label{eq:role-information-cn}
\end{equation}
$\chi_a$衡量离散变量与连续变量在同一割中的系数占比，两类变量的绝对系数和各占总量的一半时取1，仅一类参与时取0；$\mathcal H_a$
衡量五类变量系数占比的分散程度，全部系数属于同一类时为0，五类占比相同时为1。这些特征使策略能够区分通用质量相近、但变量类型及系数分布不同的候选割。

\paragraph{学习排序与候选池构建}
强化学习策略将23维特征作为候选表示，由上层确定选割数量，下层生成有序子集，并以束搜索保留多个候选序列。设$k$为
经求解器上限截断后的目标选割数。将策略序列与质量次序合并去重，并截取至多$2k$个候选形成扩展池$\mathcal E$。质
量原始分数依次按0.36、0.24、0.15、0.15、0.10和$-0.05$加权有效性、违约程度、整数支持、非其他类占比、主导已识
别角色占比和角色熵；前三项在候选内做极差归一化，再将合成分数归一化为$q_a\in[0,1]$，零极差取零。策略序列优先进
入候选池，质量排序补充其余候选。

\paragraph{基于边际增益的结构感知重排}
特征提取的结果继续用于重排：角色占比决定集合覆盖，角色集中程度参与割质量评价，通用特征与角色组成共同确定候选相
似度。令$\pi_a$为按上述完整合并排序线性递减的先验分数，$d_g=\sum_{a\in\mathcal E}p_{ag}/\sum_{a\in\mathcal E,g}p_{ag}$
为候选池中各类变量的系数占比。$s_{ab}$为候选前五维通用特征与前四类角色占比拼接向量的余弦相似度，截断至$[0,1]$且对角元
素为1。对所选集合$A\subseteq\mathcal E$，采用：
\begin{equation}
\begin{aligned}
G(A)={}&\sum_{a\in A}(0.35q_a+0.20\pi_a)\\
&+0.25\sum_gd_g\sqrt{\sum_{a\in A}p_{ag}}\\
&+\frac{0.20}{|\mathcal E|}\sum_{b\in\mathcal E}\max_{a\in A}s_{ba}.
\end{aligned}
\label{eq:cut-rerank-cn}
\end{equation}
第一项结合割质量与策略偏好；第二项以非减凹函数累计角色覆盖，同一角色被充分覆盖后，继续增加该类割的收益递减；第三
项采用设施选址形式的代表性评价，使已选割覆盖候选池中的不同特征方向。这三项均为单调次模函数，其非负加权和保持单调
次模性\citep{krause2014submodular}，据此采用边际增益贪心准则构造割子集。令$G(\varnothing)=0$，先将策略给出的
前$\lceil0.25k\rceil$个候选加入已选集合，数量不超过实际策略输出数；随后每次加入使$G(A\cup\{a\})-G(A)$最大的候选，
直至达到目标选割数量。

结构感知重排联合评价候选割的个体质量及其相对于已选集合的边际增益。角色覆盖项通过递减边际收益衡量变量角色覆盖，代
表性项降低相似候选重复入选的增量贡献，二者与策略先验共同确定割子集的选择次序。

选割策略沿用HEM的分层策略梯度框架进行离线训练\citep{wang2023hem}。训练集包含5个规模为4--12片产品的既有调度MILP实例，
每轮随机抽取实例进行8次求解采样，每次主求解限时30~s。训练时，上层随机采样选割比例，下层采样有序子集，以主求解阶段
原始--对偶积分的负值作为回报，采用带指数移动平均基线的REINFORCE更新两级策略，基线系数为0.9，上层每两轮更新一次。
共训练60轮，固定采用第60轮网络参数接入本文模型；当前配置D的实验采用这些冻结策略进行迁移评价，不在测试实例上更新网络；推理时参数保持不变，采用数量预测、束搜索和结构重排完成选割。

实现中每次学习推理至多考察128个候选、选择16个割，网络隐藏维数为128，束搜索宽度为3，并保留求解器对强制割和可选数量
的限制。为区分完整算法与组件，将RL-DMCT的学习选割组件记为RL-DMCT-S；该组件与SPBS初始解策略组合构成RL-DMCT的主求
解流程。消融方法及比较条件统一在第~\ref{sec:experiment-settings-cn}节说明。

\section{RL-DMCT的两阶段调度优化}

\paragraph{第一阶段：联合调整配对、分腔、周期与资源次序}
以通过验证的初始化结果为当前可行解，在第三章的完整可行域$\mathcal F$中按式~\eqref{eq:primary-objective-cn}最小
化$C_{\max}$。搜索同时调整产品配对、允许加工腔室内的分配、周期占位、清洁位置及搬运次序，并重新协调动作时刻。各腔室的生产
段共用VTR、装载锁与辅助片，资源约束决定各段之间的搬运和借还次序；满足工艺和驻留条件的搬运与加工可以重叠执行。每获得产品回片完工时间更
短的可行排程，即更新当前最好解。第一阶段由SCIP执行整数搜索\citep{achterberg2009scip,bestuzheva2023scip}，RL-DMCT
在选割环节提取结构特征、调用训练后的选割网络并重排候选割，随后继续联合优化排程变量。停止时，若存在通过验证的可行解，
将其赋值及产品完工时间记为$\widehat\xi$和$\widehat C$，记录有效下界、间隙与停止状态，并在启用精修时进入第二阶段；
否则记录在给定预算内未获得可行解，并跳过依赖可行输入的精修。

\paragraph{第二阶段：保留主要结构并精修时序}
以$\widehat\xi$为起点，按式~\eqref{eq:secondary-domain-cn}固定配对、分腔、周期占位、辅助身份、手分配及资源次序，
并将产品完工时间限定在$\widehat C+\varepsilon_C$内。调整动作开始时刻和取放同步选择，以式~\eqref{eq:secondary-objective-cn}
协调总等待、最大单项等待与同链相邻周期开始间隔的变化。供回片动作与同链周期的开始时刻在共享资源和驻留约束下联合调整。
通过完整验证的改善结果替换第一阶段排程，随后解码并输出晶圆动作表。

\section{批次任务解码与晶圆排程输出}
\label{sec:pjob-decoding-cn}

$\xi^\star$表示上述构造、搜索及时序精修最终保留的可行赋值；$\mathcal S^\star$是对应的具体晶圆排程。$\mathsf{Dec}$是
排程解码函数，按$\mathcal S^\star=\mathsf{Dec}(\xi^\star)$恢复实际晶圆身份和全部事件时刻。对任务$j$，将满足
$\sum_\ell x_{j\ell pr}=1$的对内位置按首次LP取片开始时刻非递减排列；同刻依次按来源LP、对编号和对内位置升序排列。
$i_j^{(k)}$是排序后的对内位置索引，表示任务$j$的第$k$片产品对应的位置。依据任务内编号映射恢复：
\begin{equation}
i_j^{(k)}\longmapsto\mathrm{Pjob}_j[k]
\longmapsto\mathcal L_j[k],\qquad k=1,\ldots,N_j.
\label{eq:pjob-decoding-cn}
\end{equation}
库存与出片顺序约束保证映射前后的来源LP一致，同刻排序规则保证映射的确定性。将恢复的实际编号关联至相应位置的全部事件，
生成具体晶圆动作表。在第~\ref{sec:model-encoding-cn}节的编码条件下，该映射只恢复身份，保留腔室、资源和动作时刻，
因而不改变产品完工时间、等待与节拍指标。

以表~\ref{tab:pjob-code-example-cn}的输入为例，若任务1的位置依次排为$(1,0),(1,1),(2,0),(3,0)$，则恢复实际编号1、3、26、27；位置$(2,0)$对应$\mathrm{Pjob}_1[3]$和实际晶圆26。前两位置同刻从LP1取片时，按对内位置分别赋予编号1和3，不改变时刻。任务2的编号2独立恢复，任务间可以交错取片。

解码后，按事件时刻生成每片晶圆的来源、PW对、腔室、处理周期及全部取放记录，并以同一时间轴独立重放资源占用。验证须确认
每片产品恰好完成规定路线与处理次数，辅助片借还区间不冲突，机械手、装载锁和腔室容量不超限，且出片顺序与驻留条件均成立。
验证通过后输出晶圆动作表、产品完工时间、全部动作完工时间及等待与节拍指标，同时保留第一阶段下界、间隙与停止状态。


\begin{table}[!htbp]
\raggedright\footnotesize
\algorithmtitle{RL-DMCT调度与晶圆排程解码}\label{alg:model4-cn}
\begin{tabularx}{\linewidth}{X}\toprule
\textbf{输入：}批次任务、LP库存与编号映射；设备参数；初始化及两阶段预算。\\
\textbf{输出：}经验证的晶圆动作表、各阶段目标与第一阶段上下界；或未取得可行解的停止记录。\\\midrule
\textbf{1. 构造并映射。}按任务资格分腔、组成PW对，安排连续加工及清洁，计算供回片时刻，验证后补齐完整模型赋值；已有有效输入可直接使用，构造失败时记录原因并转为冷启动。\\
\textbf{2. 改善起点。}启用SPBS时逐层开放资源次序与其他规定变量，只接收完整验证通过的可行赋值。\\
\textbf{3. 优化产品完工。}解除初始化临时固定关系，在完整候选域内联合调整配对、分腔、周期与资源次序；按实验臂使用原生或学习选割，保留最好可行解及有效下界。\\
\textbf{4. 精修时序。}存在有效主解且启用精修时，保持产品完工上界和主要离散结构，优化总等待、最大单项等待与节拍；无有效改善时保留主解。\\
\textbf{5. 解码与输出。}按取片顺序恢复实际编号，检查完整赋值并重放物理事件，输出可行排程、产品完工和全部占用结束时刻，以及真实停止状态。\\\bottomrule
\end{tabularx}
\end{table}
\FloatBarrier
\clearpage
\section{产品完工瓶颈驱动的原始解改进方法}\label{sec:third-innovation}
\subsection{瓶颈识别与候选调整对象}
\vspace*{5cm}
\subsection{优质排程结构共识与分歧}
\vfill\null\clearpage
\subsection{联合调整邻域与MILP修复}
\par\vspace*{4cm}
\subsection{质量引导搜索、预算分配与终止}
\vspace*{3cm}
\subsection{可行性保持及最优性判据}
\vfill\null
